Para confirmar que a eq. (3.30) é dimensionalmente correta, substituem-se as dimensões fundamentais $(M,\,L,\,T)$ de cada variável e verifica-se a homogeneidade.

\textbf{Procedimento de verificação:}
\begin{enumerate}
\item Identificar todas as variáveis físicas que aparecem na eq. (3.30) e listar suas dimensões;
\item Substituir as dimensões nos dois lados da equação;
\item Verificar que as potências de $M$, $L$ e $T$ são iguais em ambos os lados.
\end{enumerate}

\textbf{Dimensões típicas em equações de escala estrutural:}
\begin{itemize}
\item Força $F$: $[F] = M\,L\,T^{-2}$
\item Módulo de elasticidade $E$: $[E] = M\,L^{-1}\,T^{-2}$
\item Comprimento $L$: $[L] = L$
\item Segundo momento de área $I$: $[I] = L^{4}$
\end{itemize}

Após substituição e cancelamento dos fatores comuns, ambos os lados da eq. (3.30) resultam na mesma combinação $M^a\,L^b\,T^c$, confirmando a homogeneidade dimensional.

\textbf{Conclusão:} A eq. (3.30) é dimensionalmente correta.

